% Fragmento público generado desde propietarios íntegros.
% Residencia: 52.13.2.
% Fuente: ../incoming/radion_monografia_20260827/manuscrito/sections/13_sintesis.tex:135-184
\subsubsection{Système de théorèmes composants du radion}

\begin{theorem}[Généalogie]
L'objet \(\mathfrak R_\sigma\) descend de
\(\APP\to\TRIT\to\TPK\to X_{\mathrm{enr}}\) et ne reçoit pas la cible métrologique
comme sélecteur.
\end{theorem}

\begin{theorem}[Charnière]
Dans la carte APP de cent marques, la couture critique et la deuxième phase
dodécaphasique satisfont \(j_{100}(1/2)=\theta_2=50^\circ\), avec \(m=2\) unique.
\end{theorem}

\begin{theorem}[Densité]
Les deux feuillets du secteur actif \(N_6=90\), normalisés par \(3^6=729\),
produisent \(\rho_{\mathrm{tw}}=20/81\).
\end{theorem}

\begin{theorem}[Retenue]
La soustraction APP des sections \(S_T\) et \(S_E\) converge vers
\(\epsone\) avec un taux d'intervalle \(3^{-54}\) par retour et ferme exactement
l'unité.
\end{theorem}

\begin{theorem}[Action]
L'ellipse positive satisfait \(\mathcal A(\theta)=\hret\theta\) ; un radion
balaie \(\hret\) et un tour enferme \(\hfull\).
\end{theorem}

\begin{theorem}[Bord--intérieur]
La même frontière possède une action finie \(\hfull\), tandis que l'intérieur de
Klein a une aire hyperbolique infinie.
\end{theorem}

\begin{theorem}[Forme]
Le quotient \(e^{-\eta}\) est conservé à travers les demi-axes spectraux,
les demi-axes d'action, les incertitudes, l'impédance et le module du tore.
\end{theorem}

\begin{theorem}[Polarité]
L'inversion \(C^*\mapsto-C^*\) échange les lois constitutives, inverse
l'impédance et conserve la vitesse normalisée.
\end{theorem}

\begin{theorem}[Modularité]
L'inversion bidirectionnelle agit comme \(S:\tau\mapsto-1/\tau\) et conserve
\(j(\tau)\), tandis que l'étiquetage focal conserve l'orientation que
\(j\) oublie.
\end{theorem}

% Fuente: ../incoming/radion_monografia_20260827/manuscrito/sections/A_demostraciones.tex:1-179
\subsubsection{Démonstrations algébriques et géométriques réunies}

\paragraph{Démonstration linéaire de la clôture unitaire}

Pour faciliter un audit indépendant, nous réunissons l'argument sans
interruption narrative.

\begin{enumerate}[label=\textbf{Étape \arabic*.},leftmargin=2.2cm]
\item La connexion nonadique génère de manière corrélée
\(\pih,e_{\HMT},\varphi_{\HMT},\alphah\).
\item Le tour est \(T_+=2\pih\).
\item La roue fixe \(\theta_2=50^\circ\), et la carte parangulaire fixe
\(A=1000\alphah\).
\item La section directe est
\[
S_E=T_+\left(\frac{50}{360}+\frac{1000\alphah}{360}\right)
=2\pih\left(\frac5{36}+\frac{25}{9}\alphah\right).
\]
\item Les intervalles certifiés donnent \(1<S_E<2\) ; donc
\(D_E=S_E-1\).
\item L'incidence produit \(N_6=90\) ; deux feuillets et l'espace ambiant \(729\)
produisent \(\rho=2N_6/729=20/81\).
\item La torsion globale est
\[
\Delta_4=\frac{\pih-e_{\HMT}\log\pih}{270}
\left(1+\frac{\pih}{729}\right).
\]
\item La section torsionnelle est \(S_T=1+(20/81)\Delta_4\).
\item La soustraction APP définit
\[
\epsone=S_T-S_E=\frac{20}{81}\Delta_4-D_E.
\]
\item Par substitution,
\[
S_E-\frac{20}{81}\Delta_4+\epsone=1.
\]
\item Multiplier par \(360/T_+=180/\pih\) produit
\[
\frac{180^\circ}{\pih}
=50^\circ+A^\circ-\frac{20}{81}\Delta_4^\circ+\epsR^\circ.
\]
\end{enumerate}

Aucune étape ne reçoit le membre de gauche de la dernière égalité. Le résultat
métrologique n'est évalué qu'à la fin comme reconnaissance.

\paragraph{Unicité de la soustraction avec retenue}

Soit \(B=1000\). Pour chaque position, tout entier \(u_i\) admet une division
euclidienne unique
\[
u_i=Bc_i+r_i,\qquad 0\le r_i<B.
\]
Par conséquent,
\[
r_i=u_i\bmod B,\qquad c_i=(u_i-r_i)/B
\]
sont uniques. En parcourant les positions de l'extrémité distante vers
l'avant, la retenue de sortie d'une position fixe l'entrée de la suivante.
Par récurrence, tout le mot de restes et de retenues est unique.

En pseudocode :
\begin{verbatim}
acarreo = remote_acarreo
for i from last_block downto first_block:
    u = torsion[i] - direct[i] + acarreo
    remainder[i] = u mod 1000
    acarreo = (u - remainder[i]) / 1000
return remainder, acarreo_word
\end{verbatim}
La procédure ne consulte pas la valeur attendue de la différence.

\paragraph{Convergence de la fonctionnelle selon la profondeur}

Soit
\[
F(p,e)=1+\frac{20}{81}
\frac{p-e\log p}{270}\left(1+\frac{p}{729}\right)
-2p\left(\frac5{36}+\frac{25}{9}\alphah\right).
\]
Dans le rectangle compact \(3\le p\le4\), \(2\le e\le3\), \(F\) est
continue et continûment différentiable. Si \(I_N^\pi\searrow\{\pih\}\) et
\(I_N^e\searrow\{e_{\HMT}\}\), l'image par intervalles
\[
F(I_N^\pi,I_N^e)
\]
forme une famille emboîtée dont la largeur tend vers zéro. L'intersection est un
unique réel, \(\epsone\). La monotonie des dérivées permet d'évaluer les
extrémités sans échantillonnage intérieur.

Le taux \(3^{-54}\) par retour dérive du raffinement de six trits sur
neuf niveaux. À la profondeur de \(45\) triades, la borne de largeur
\(<4.81\times10^{-130}\) certifie plus de 129 chiffres de la valeur.

\paragraph{Dérivation de la métrique et de l'aire hyperbolique de Klein}

Dans le disque unité, la forme matricielle de la métrique de Klein est
\[
g(u,v)=\frac1{1-r^2}I
+\frac1{(1-r^2)^2}
\begin{pmatrix}u\\v\end{pmatrix}
\begin{pmatrix}u&v\end{pmatrix}.
\]
Le lemme du déterminant pour une perturbation de rang un donne
\[
\det g=(1-r^2)^{-3}.
\]
D'où
\[
\sqrt{\det g}\,\diff u\diff v
=(1-r^2)^{-3/2}\diff u\diff v.
\]
L'intégration en coordonnées polaires produit \eqref{eq:klein-area-R}. La limite infinie est
due à la singularité métrique à la frontière ; la frontière conserve une
aire symplectique finie parce que cette mesure utilise \(\diff Q\wedge\diff P\),
non \(\sqrt{\det g}\).

\paragraph{Distance focale dans le disque de Klein}

Sur un diamètre, la distance du centre à \(r\) est
\[
d_K(0,r)=\operatorname{artanh}r.
\]
Pour \(r=e_f\), \eqref{eq:excentricidad} implique
\[
1-e_f^2=e^{-2\eta}.
\]
Si \(u_f=\operatorname{artanh}e_f\), alors
\[
\operatorname{sech}u_f=\sqrt{1-e_f^2}=e^{-\eta},
\]
et donc \(\eta=\log\cosh u_f\), comme dans
\eqref{eq:focal-hyperbolic-chain}.

\paragraph{Transformation symplectique de compression}

Définissons
\begin{equation}
S_\eta=
\begin{pmatrix}e^{\eta/2}&0\\0&e^{-\eta/2}\end{pmatrix},
\qquad \det S_\eta=1.
\label{eq:compresión simpléctica}
\end{equation}
La structure complexe transportée est
\begin{equation}
J_\eta=S_\eta
\begin{pmatrix}0&-1\\1&0\end{pmatrix}
S_\eta^{-1}
=\begin{pmatrix}0&-e^\eta\\e^{-\eta}&0\end{pmatrix},
\qquad J_\eta^2=-I.
\label{eq:Jeta}
\end{equation}
La courbe
\[
\gamma(\theta)=\sqrt{2\hret}\,
S_\eta(\cos\theta,\sin\theta)
=e^{\theta J_\eta}\gamma(0)
\]
est exactement l'ellipse d'action. La compression symplectique hyperbolique conserve l'aire et
transporte la rotation complexe au lieu de la détruire.

Pour l'orientation \(\sigma\),
\[
\gamma_\sigma(\theta)
=(a_S\cos\theta,\sigma b_S\sin\theta),
\qquad
\mathcal A_\sigma(\theta)=\sigma\hret\theta.
\]
Les deux feuillets partagent les foyers et le module d'action ; ils diffèrent par l'action signée.

\paragraph{Vérification diophantienne du poids dodécaphasique de Barbero--Immirzi}

L'équation \(4a-3b=45\) a pour solutions entières
\[
(a,b)=(12+3t,1+4t),\qquad t\in\Z.
\]
La condition \(|b|=1\) laisse \(b=1\), car \(b=-1\) n'appartient pas à la classe
\(1\pmod4\). D'où \(a=12\). Cette vérification reproduit le couple obtenu par la chaîne fondamentale
dodécaphasique : douze secteurs et une unique couture orientée. Les cardinaux
\(90/120\), les multiplicités de la chaîne et le lecteur de retour
conservent des fonctions distinctes ; le coefficient de pôle et la minimalité
vérifient arithmétiquement leur composition.
